\pagestyle{empty}

\section*{A mellékelt adathordozó tartalma}
% \addcontentsline{toc}{chapter}{A mellékelt adathordozó tartalma}

Az adathordozó a következőket tartalmazza:

\begin{itemize}
\item \texttt{dolgozat.pdf}: A dolgozat PDF formátumban.
\item \texttt{hasznalati.pdf}: A használati útmutató PDF formátumban.
\item \texttt{dolgozat/}: A dolgozat LaTeX forráskódja. Fordításához szükségesek a \texttt{pdflatex} és \texttt{bibtex} programok. Az \texttt{svg} csomag dependenciájaként az \texttt{Inkscape} program is szükséges. 
\item \texttt{program/}: A szerkesztő teljes forráskódja a tesztekkel együtt. Az \texttt{src} mappa az alkalmazás kódját, a \texttt{tests} mappa az automatikus teszteket, a \texttt{docs} mappa az architektúra leírását, a tervezési döntéseket és egy angol nyelvű felhasználói útmutatót tartalmaz. A \texttt{dist} mappában az alkalmazás lefordított, éles változata, az \texttt{eval} mappában pedig \aref{sec:measurements}. szakasz mérőszkriptjei és a mérésekhez felhasznált fájlok listája található.
\end{itemize}

\subsection*{A szerkesztő futtatása telepítés nélkül}

A \texttt{program/dist} mappa tartalmát bármely statikus webszerver kiszolgálhatja. Ha a gépen van Python 3, a mappában kiadott \texttt{python3 -m http.server 8000} parancs után a szerkesztő a \texttt{http://localhost:8000} címen nyitható meg. Az \texttt{index.html} fájl közvetlen megnyitása nem elég, mert a böngészők a \texttt{file://} címről betöltött JavaScript-modulokat biztonsági okból általában letiltják.

\subsection*{Fordítás a forráskódból}

A fordításhoz Node.js 20.19 vagy 22.12, illetve ezeknél újabb változat szükséges, ezt a Vite 7 írja elő. A \texttt{program} mappában kiadandó parancsok:

\begin{description}
\item[\texttt{npm install}] telepíti a függőségeket, ehhez internetkapcsolat szükséges;
\item[\texttt{npm run dev}] elindítja a fejlesztői szervert, a szerkesztő a kiírt címen, alapértelmezés szerint a \texttt{http://localhost:5173} címen nyitható meg böngészőben;
\item[\texttt{npm run build}] újra elkészíti az éles változatot a \texttt{dist} mappába;
\item[\texttt{npm run preview}] helyben kiszolgálja az elkészült éles változatot;
\item[\texttt{npm test}] lefuttatja az automatikus teszteket.
\end{description}

A szerkesztő egy Chromium-alapú böngészőben, például Google Chrome vagy Microsoft Edge alatt használható a legteljesebb funkcionalitással, mert a helyi betűtípusok elérése csak ezekben támogatott. A Google Fonts-betűtípusok letöltéséhez internetkapcsolat szükséges, a szerkesztő többi része hálózat nélkül is működik.
